\section{Negative results, with exact scope}\label{sec:negative}

The record contains several obstructions: proved limits of particular methods, and sequences shown not to give what was sought. Each is stated here with its proof or its source, and with what it does not show. Each concerns a method, a sequence or an observable, and several of them name what the method would need.

\subsection{The global perturbative route}

\begin{proposition}[the collapse of the global analyticity disk; spatial line §27]\label{prop:neumann}
Write $H=(2g^2/a)(H_0-\xi W_L)+\mathrm{const}$ on the full space (no gauge projection), with $H_0=\sum_eE_e$ the sum of the link Casimirs $E_e$ and $W_L=\sum_pW_p$. The operator-norm Neumann series certifies the isolated rank-one ground-state projection on the circle $|z|=3/8$ exactly for $|\xi|<3/(16M_L)=1/(64L^2(2L+1))$, and on no contour separating $0$ from the rest of the spectrum does the norm-product criterion give more. On the physical space the same route gives exactly $|\xi|<3/(4M_L)$, on the circle $|z|=3/2$.
\end{proposition}
\status{Proved here; the full-space statement audited (YM-07 §27); the physical-space radius and the optimality over contours are \textsf{Generalised here}.}
\begin{proof}
$|W_p|\le2$ with equality at $U\equiv I$, so $\|\xi W_L\|=2|\xi|M_L$; the same holds on the physical space, since for every $\varepsilon>0$ the indicator of the nonempty open gauge-invariant set $\{W_L>2M_L-\varepsilon\}$ is a physical vector. The spectrum of $H_0$ is contained in $\{0\}\cup[3/4,\infty)$, so $\|(H_0-z)^{-1}\|\le8/3$ on $|z|=3/8$, and the series converges when $2|\xi|M_L\cdot\frac83<1$. Both factors are attained ($\|(H_0-z)^{-1}\|=8/3$ at $z=\pm3/8$), so the norm-product criterion fails beyond this radius. Any contour separating $0$ from the rest of the spectrum crosses $(0,\frac34)$, since the eigenvalue $\frac34$ must have winding number $0$; at a crossing point $x$ the resolvent norm is $\max\{1/x,1/(\frac34-x)\}\ge\frac83$. On the physical space the spectrum of $H_0$ lies in $\{0\}\cup[3,\infty)$ by Lemma~\ref{lem:casimir}, with $3$ attained by one plaquette in spin $\frac12$, and the same two arguments give the resolvent norm $\frac23$ on $|z|=\frac32$ and the radius $3/(4M_L)$.
\end{proof}
The radius shrinks with the box, so at any fixed coupling this route eventually fails; it does not show that the true analyticity radius is small, and it proves neither a gap nor its absence. It explains why Theorem~\ref{thm:ym01} needs local norms: at $L=2$ ($M_2=240$) the global route certifies only $|\xi|<1/1280$ on the full space and $|\xi|<1/320$ on the physical space, both far inside $\xi\le\alpha_5\approx1/54.3$. (\emph{Weak-coupling disk} in the source means small $|\xi|$, that is large $g$: the regime of Theorem~\ref{thm:ym01}.)


\subsection{The other obstructions}

\begin{center}\small
\begin{longtable}{@{}>{\raggedright\arraybackslash}p{0.19\linewidth}>{\raggedright\arraybackslash}p{0.29\linewidth}>{\raggedright\arraybackslash}p{0.23\linewidth}>{\raggedright\arraybackslash}p{0.2\linewidth}@{}}
\toprule
Source & What fails & Status & What it does not show\\
\midrule
\endhead
fifth reference (§H7), and Proposition~\ref{prop:path} & the strong-coupling estimates along the standing continuum path & proved here & a statement about the continuum regime, in either direction\\
fifth reference (§H7) & re-centring the elementary (order-1, single-norm) certificate to extend the source domain: it returns its old endpoint $3/256=\alpha_1$ & located (the workbench's \texttt{ROUTE\_ASSESSMENT} calculation) & that other routes fail; the workbench names the next candidate\\
quantum line §11 & lower bounds for dressed channels as bounds for the whole spectral measure or for all physical states & the document's own statement & a bound on the gap by other means\\
volume line §26 & the certified spectral window of the weighted covariance stays bounded along the logarithmic path: its upper end diverges & located & that the covariance has no weight in a fixed window\\
spatial line §§17--18 & a selected weak-coupling diagonal as the sought interacting continuum: the finite gaps close while the magnetic probability escapes & located & the prescribed logarithmic and fixed-electric-coefficient paths, or other continuations\\
spatial line §20 & a unique-vacuum realisation of a selected macroscopic limit of a second observable ($d\delta_0$, $d=1-2^{-3/2}$) & located & anything beyond that selected limit\\
quantum line, Theorem 14.2 (Jacobian branch) & a mass-gap conclusion from the trial states: their energies are the free two-gluon threshold of a shrinking-coupling box & audited in \note{34\_} (Section~\ref{sec:side}) & a statement about the continuum mass gap\\
quantum line §§20--22 (Navier--Stokes branch) & a zero energy quotient at fixed coupling: the zero-quotient sequence exists only along couplings tending to $0$ & checked in \note{21\_} §2, conditional on the imported rate bounds & whether another state with vanishing quotient exists\\
\bottomrule
\end{longtable}
\end{center}
All of these concern methods, sequences or observables. The workbench states each of them itself, in the passages cited; in the passages read, no negative result of the record is stated more broadly than its proof.
